
\documentclass{article}
\usepackage{colortbl}
\usepackage{makecell}
\usepackage{multirow}
\usepackage{supertabular}

\begin{document}

\newcounter{utterance}

\centering \large Interaction Transcript for game `referencegame', experiment `random\_grids', episode 1 with reflexion{-}llama{-}3.1{-}8b{-}instruct.
\vspace{24pt}

{ \footnotesize  \setcounter{utterance}{1}
\setlength{\tabcolsep}{0pt}
\begin{supertabular}{c@{$\;$}|p{.15\linewidth}@{}p{.15\linewidth}p{.15\linewidth}p{.15\linewidth}p{.15\linewidth}p{.15\linewidth}}
    \# & $\;$A & \multicolumn{4}{c}{Game Master} & $\;\:$B\\
    \hline

    \theutterance \stepcounter{utterance}  
    & & \multicolumn{4}{p{0.6\linewidth}}{
        \cellcolor[rgb]{0.9,0.9,0.9}{
            \makecell[{{p{\linewidth}}}]{
                \texttt{\tiny{[P1$\langle$GM]}}
                \texttt{You have attempted this same task before and did not succeed. These are the plans you wrote after those attempts.} \\
\texttt{Trial \#1: New plan: I failed to distinguish the target grid from the distractors because my description was too focused on specific positions of X's in the grid. What I should have done is to look for a unique pattern or feature that is present in the target grid but not in the distractors. Next time, I will carefully examine the grids and try to identify a distinctive characteristic that can be used to describe the target grid. For example, I might look for a row or column that has a unique combination of X's and empty cells.} \\
\\ 
\texttt{Trial \#2: New plan: I failed to distinguish the target grid from the distractors because my description was too focused on specific positions of X's in the grid. What I should have done is to look for a unique pattern or feature that is present in the target grid but not in the distractors. Next time, I will carefully examine the grids and try to identify a distinctive characteristic that can be used to describe the target grid. For example, I might look for a row or column that has a unique combination of X's and empty cells. I will also pay closer attention to the distribution of X's in the target grid and try to find a feature that is not present in the distractors.} \\
\\ 
\texttt{Trial \#3: New plan: I failed to distinguish the target grid from the distractors because my description was too focused on specific positions of X's in the grid. What I should have done is to look for a unique pattern or feature that is present in the target grid but not in the distractors. Next time, I will carefully examine the grids and try to identify a distinctive characteristic that can be used to describe the target grid. For example, I might look for a row or column that has a unique combination of X's and empty cells. I will also pay closer attention to the distribution of X's in the target grid and try to find a feature that is not present in the distractors. Additionally, I will consider the overall structure of the grid, such as the number of rows or columns with a specific number of X's, and use this information to create a more accurate and distinctive description.} \\
\\ 
\\ 
\texttt{You are given three grids, where each of them is 5 by 5 in size.} \\
\texttt{Grids have empty cells marked with "▢" and filled cells marked with "X".} \\
\texttt{Your task is to generate a referring expression that best describes the target grid while distinguishing it from the two other distractor grids.} \\
\texttt{The first grid is the target grid, and the following two grids are the distractors.} \\
\\ 
\texttt{Target grid:} \\
\\ 
\texttt{□ X □ X □} \\
\texttt{X □ X □ X} \\
\texttt{□ X □ X □} \\
\texttt{X □ X □ X} \\
\texttt{□ X □ X □} \\
\\ 
\texttt{Distractor grid 1:} \\
\\ 
\texttt{□ X X □ X} \\
\texttt{X □ X X □} \\
\texttt{□ X □ X X} \\
\texttt{X □ X □ X} \\
\texttt{□ X X □ □} \\
\\ 
\texttt{Distractor grid 2:} \\
\\ 
\texttt{□ □ X X □} \\
\texttt{X □ □ X X} \\
\texttt{□ X X □ X} \\
\texttt{X X □ □ X} \\
\texttt{□ X X □ □} \\
\\ 
\texttt{Instruction: Describe the target grid.} \\
\texttt{Generate the referring expression starting with the tag "Expression: " for the given target grid. Omit any other text.} \\
            }
        }
    }
    & \\ \\

    \theutterance \stepcounter{utterance}  
    & \multicolumn{4}{p{0.6\linewidth}}{
        \cellcolor[rgb]{0.9,0.9,0.9}{
            \makecell[{{p{\linewidth}}}]{
                \texttt{\tiny{[P1$\rangle$GM]}}
                \texttt{Expression: The grid that has a row with two X's in the first and last columns and a row with two X's in the fourth and fifth columns.} \\
            }
        }
    }
    & & \\ \\

    \theutterance \stepcounter{utterance}  
    & & & \multicolumn{2}{p{0.3\linewidth}}{
        \cellcolor[rgb]{0.9,0.9,0.9}{
            \makecell[{{p{\linewidth}}}]{
                \texttt{\tiny{[GM$|$GM]}}
                \texttt{Expression: The grid that has a row with two X's in the first and last columns and a row with two X's in the fourth and fifth columns.} \\
            }
        }
    }
    & & \\ \\

    \theutterance \stepcounter{utterance}  
    & & \multicolumn{4}{p{0.6\linewidth}}{
        \cellcolor[rgb]{0.9,0.9,0.9}{
            \makecell[{{p{\linewidth}}}]{
                \texttt{\tiny{[GM$\rangle$P2]}}
                \texttt{You have attempted this same task before and did not succeed. These are the plans you wrote after those attempts.} \\
\texttt{Trial \#1: New plan: I failed to match the grid with the given expression because I didn't carefully analyze the expression. I should have broken down the expression into its individual components and checked each grid against them. Specifically, I should have checked the first row of each grid for exactly two X's, the last row for exactly two X's, and the third column for exactly two X's. I should have also checked the grids more systematically, rather than just selecting the first one. Next time, I will carefully analyze the expression and check each grid against its components before making a selection.} \\
\\ 
\texttt{Trial \#2: New plan: I failed to match the grid with the given expression because I didn't carefully analyze the expression. I should have broken down the expression into its individual components and checked each grid against them. Specifically, I should have checked the first row of each grid for exactly two X's, the last row for exactly two X's, and the third column for exactly two X's. However, upon re{-}examining the expression, I realized that it also mentions a row with two X's in the fourth and fifth columns. I should have checked the second row of each grid for exactly two X's in the fourth and fifth columns. Next time, I will carefully analyze the expression and check each grid against its components, including the row with two X's in the fourth and fifth columns, before making a selection.} \\
\\ 
\texttt{Trial \#3: New plan: I failed to match the grid with the given expression because I didn't carefully analyze the expression. I should have broken down the expression into its individual components and checked each grid against them. Specifically, I should have checked the first row of each grid for exactly two X's, the last row for exactly two X's, the third column for exactly two X's, and the second row for exactly two X's in the fourth and fifth columns. However, upon re{-}examining the expression, I realized that it also mentions a column with two X's and three empty cells in the second column. I should have checked the second column of each grid for exactly two X's and three empty cells. Next time, I will carefully analyze the expression and check each grid against its components, including the column with two X's and three empty cells in the second column, before making a selection.} \\
\\ 
\\ 
\texttt{You are given three grids, where each of them is 5 by 5 in size.} \\
\texttt{Grids have empty cells marked with "▢" and filled cells marked with "X".} \\
\texttt{You are also given a referring expression that describes one of the given grids.} \\
\texttt{Your task is to select the grid that matches the given referring expression.} \\
\texttt{Generate only the number (in text) of the grid that matches the given expression by selecting first, second, or third.} \\
\\ 
\texttt{First:} \\
\\ 
\texttt{□ X X □ X} \\
\texttt{X □ X X □} \\
\texttt{□ X □ X X} \\
\texttt{X □ X □ X} \\
\texttt{□ X X □ □} \\
\\ 
\texttt{Second:} \\
\\ 
\texttt{□ X □ X □} \\
\texttt{X □ X □ X} \\
\texttt{□ X □ X □} \\
\texttt{X □ X □ X} \\
\texttt{□ X □ X □} \\
\\ 
\texttt{Third:} \\
\\ 
\texttt{□ □ X X □} \\
\texttt{X □ □ X X} \\
\texttt{□ X X □ X} \\
\texttt{X X □ □ X} \\
\texttt{□ X X □ □} \\
\\ 
\texttt{Expression: The grid that has a row with two X's in the first and last columns and a row with two X's in the fourth and fifth columns.} \\
\texttt{Question: Which grid does the expression refer to?} \\
\texttt{Start with the tag "Answer: ", followed by your selection. Omit any other text.} \\
            }
        }
    }
    & \\ \\

    \theutterance \stepcounter{utterance}  
    & & & \multicolumn{4}{p{0.6\linewidth}}{
        \cellcolor[rgb]{0.9,0.9,0.9}{
            \makecell[{{p{\linewidth}}}]{
                \texttt{\tiny{[GM$\langle$P2]}}
                \texttt{Answer: 1} \\
            }
        }
    }
     \\ \\

    \theutterance \stepcounter{utterance}  
    & & & \multicolumn{2}{p{0.3\linewidth}}{
        \cellcolor[rgb]{0.9,0.9,0.9}{
            \makecell[{{p{\linewidth}}}]{
                \texttt{\tiny{[GM$|$GM]}}
                \texttt{1} \\
            }
        }
    }
    & & \\ \\

    \theutterance \stepcounter{utterance}  
    & \multicolumn{4}{p{0.6\linewidth}}{
        \cellcolor[rgb]{0.95,0.95,0.95}{
            \makecell[{{p{\linewidth}}}]{
                \texttt{\tiny{[P2$\rightarrow$P1]}}
                \texttt{Post{-}episode Reflexion (not sent to GM):} \\
\texttt{New plan: I failed to distinguish the target grid from the distractors because my description was too focused on specific positions of X's in the grid. What I should have done is to look for a unique pattern or feature that is present in the target grid but not in the distractors. Next time, I will carefully examine the grids and try to identify a distinctive characteristic that can be used to describe the target grid. For example, I might look for a row or column that has a unique combination of X's and empty cells. I will also pay closer attention to the distribution of X's in the target grid and try to find a feature that is not present in the distractors. Additionally, I will consider the overall structure of the grid, such as the number of rows or columns with a specific number of X's, and use this information to create a more accurate and distinctive description. Furthermore, I will also look for any unique relationships between the X's and empty cells in the target grid, such as a specific pattern of X's surrounding an empty cell, and use this information to create a more detailed and accurate description.} \\
            }
        }
    }
    & & \\ \\

    \theutterance \stepcounter{utterance}  
    & & & \multicolumn{4}{p{0.6\linewidth}}{
        \cellcolor[rgb]{0.95,0.95,0.95}{
            \makecell[{{p{\linewidth}}}]{
                \texttt{\tiny{[P1$\rightarrow$P2]}}
                \texttt{Post{-}episode Reflexion (not sent to GM):} \\
\texttt{New plan: I failed to match the grid with the given expression because I didn't carefully analyze the expression. I should have broken down the expression into its individual components and checked each grid against them. Specifically, I should have checked the first row of each grid for exactly two X's, the last row for exactly two X's, the third column for exactly two X's, the second row for exactly two X's in the fourth and fifth columns, and the second column for exactly two X's and three empty cells. However, upon re{-}examining the expression, I realized that I should have also checked the grids more systematically, rather than just selecting the first one. Next time, I will carefully analyze the expression and check each grid against its components, including the column with two X's and three empty cells in the second column, and check the grids more systematically before making a selection.} \\
            }
        }
    }
     \\ \\

\end{supertabular}
}

\end{document}
